\documentclass{article}
\usepackage{amsmath,amssymb,amsthm}
\usepackage{algorithm}
\usepackage[noend]{algpseudocode}
\usepackage{enumitem}
\usepackage{hyperref}
\usepackage{xcolor}
\newtheorem{theorem}{Theorem}
\newtheorem{lemma}{Lemma}
\newtheorem{assumption}{Assumption}
\newtheorem{corollary}{Corollary}
\newcommand{\gap}{g}
\newcommand{\diam}{\operatorname{diam}}
\newcommand{\conv}{\operatorname{conv}}
\newcommand{\argmin}{\operatornamewithlimits{arg\,min}}
\newcommand{\argmax}{\operatornamewithlimits{arg\,max}}
\newcommand{\lin}{\operatorname{lin}}
\newcommand{\aff}{\operatorname{aff}}
\newcommand{\pyr}{\delta_{\mathcal{P}}}
\newcommand{\C}{\mathcal{C}}
\newcommand{\V}{\mathcal{V}}
\newcommand{\A}{\mathcal{A}}

\begin{document}

\section*{Away‑step Frank‑Wolfe (AFW)}

\section*{Mathematical Formulation}
Let $f:\mathbb{R}^d\rightarrow\mathbb{R}$ be a convex and $L$‑smooth function and let
\[
\C \subset \mathbb{R}^d
\]
be a compact polytope, i.e. $\C=\operatorname{conv}(\V)$ for a finite set of vertices $\V=\{v_1,\dots,v_m\}$.
Given the current iterate $x_k\in\C$, define the \emph{Frank‑Wolfe (FW) vertex}
\[
s_k \;:=\; \argmin_{s\in\C}\;\langle \nabla f(x_k),\,s\rangle,
\tag{1}
\]
and the \emph{away vertex}
\[
v_k \;:=\; \argmax_{v\in\A_k}\;\langle \nabla f(x_k),\,v\rangle,
\tag{2}
\]
where $\A_k\subseteq \V$ is the current \emph{active set} (the set of vertices with positive coefficient in the representation of $x_k$).  

Define the FW direction $d^{\text{FW}}_k:=s_k-x_k$ and the away direction $d^{\text{A}}_k:=x_k-v_k$.  
Choose the \emph{descent direction}
\[
d_k \;:=\;
\begin{cases}
d^{\text{FW}}_k & \text{if }\langle \nabla f(x_k),d^{\text{FW}}_k\rangle\le \langle \nabla f(x_k),d^{\text{A}}_k\rangle,\\[4pt]
d^{\text{A}}_k   & \text{otherwise}.
\end{cases}
\tag{3}
\]

The maximal feasible step size for an away step is
\[
\gamma_{\max}:=
\begin{cases}
1 & \text{if }d_k=d^{\text{FW}}_k,\\[4pt]
\displaystyle\frac{\alpha_{v_k}}{1-\alpha_{v_k}} & \text{if }d_k=d^{\text{A}}_k,
\end{cases}
\tag{4}
\]
where $\alpha_{v_k}>0$ is the coefficient of $v_k$ in the current convex combination of $x_k$.

A line‑search (or exact quadratic line‑search for $L$‑smooth $f$) yields
\[
\gamma_k \;:=\;\argmin_{\gamma\in[0,\gamma_{\max}]}\; f\bigl(x_k+\gamma d_k\bigr).
\tag{5}
\]

Finally the update is
\[
x_{k+1}=x_k+\gamma_k d_k.
\tag{6}
\]

\section*{Pseudocode}
\begin{algorithm}[H]
\caption{Away‑step Frank‑Wolfe (AFW)}\label{alg:afw}
\begin{algorithmic}[1]
\Require Convex $L$‑smooth $f$, polytope $\C=\conv(\V)$, tolerance $\varepsilon>0$, max‑iter $T$.
\State Choose $x_0\in\C$, initialise active set $\A_0:=\{v\in\V: x_0=\sum_{v\in\A_0}\alpha_v v,\;\alpha_v>0\}$.
\For{$k=0,\dots,T-1$}
    \State Compute gradient $g_k:=\nabla f(x_k)$.
    \State \textbf{FW vertex:} $s_k:=\argmin_{s\in\C}\langle g_k,s\rangle$.
    \State \textbf{Away vertex:} $v_k:=\argmax_{v\in\A_k}\langle g_k,v\rangle$.
    \State Compute FW and away directions:
          $d^{\text{FW}}_k:=s_k-x_k$, $d^{\text{A}}_k:=x_k-v_k$.
    \If{$\langle g_k,d^{\text{FW}}_k\rangle\le \langle g_k,d^{\text{A}}_k\rangle$}
        \State $d_k\gets d^{\text{FW}}_k$, $\gamma_{\max}\gets 1$.
    \Else
        \State $d_k\gets d^{\text{A}}_k$, $\gamma_{\max}\gets \frac{\alpha_{v_k}}{1-\alpha_{v_k}}$.
    \EndIf
    \State Perform line‑search:
          $\gamma_k:=\argmin_{\gamma\in[0,\gamma_{\max}]} f(x_k+\gamma d_k)$.
    \State Update $x_{k+1}\gets x_k+\gamma_k d_k$.
    \State Update active set $\A_{k+1}$ (add $s_k$ if FW step, remove $v_k$ if away step finishes).
    \State Compute duality gap $\gap_k:=\langle g_k, x_k-s_k\rangle$.
    \If{$\gap_k\le\varepsilon$}
        \State \textbf{break}
    \EndIf
\EndFor
\State \Return $x_{k+1}$.
\end{algorithmic}
\end{algorithm}

\section*{Dry Run Example}
Consider the univariate quadratic
\[
f(w)=(w-3)^2,\qquad \C=[0,5]=\conv\{0,5\}.
\]
The gradient is $\nabla f(w)=2(w-3)$.  
We initialise $x_0=0$ (so $\A_0=\{0\}$) and use a fixed step size $\eta=0.1$ for illustration (the line‑search for a quadratic yields the exact minimiser, but we keep $\eta$ to mimic a back‑tracking rule).

\begin{enumerate}[label=\textbf{Iter \arabic*:}, leftmargin=*]
\item \textbf{Iteration 0}
\begin{align*}
g_0&=2(x_0-3)= -6,\\
s_0&=\argmin_{s\in\{0,5\}}\langle g_0,s\rangle =\argmin\{-6s\}=5,\\
v_0&=\argmax_{v\in\{0\}}\langle g_0,v\rangle =0,\\
d^{\text{FW}}_0&=5-0=5,\qquad d^{\text{A}}_0=0-0=0.
\end{align*}
Since $\langle g_0,d^{\text{FW}}_0\rangle=-30 < \langle g_0,d^{\text{A}}_0\rangle=0$, we take a FW step:
\[
\gamma_0=\eta=0.1,\qquad x_1=0+0.1\cdot5=0.5.
\]
Objective: $f(x_1)=(0.5-3)^2=6.25$.

\item \textbf{Iteration 1}
\begin{align*}
g_1&=2(0.5-3)=-5,\\
s_1&=\argmin_{s\in\{0,5\}}(-5s)=5,\\
v_1&=\argmax_{v\in\{0,5\}}\{-5v\}=0,\\
d^{\text{FW}}_1&=5-0.5=4.5,\qquad d^{\text{A}}_1=0.5-0=0.5.
\end{align*}
\[
\langle g_1,d^{\text{FW}}_1\rangle=-5\cdot4.5=-22.5,\qquad 
\langle g_1,d^{\text{A}}_1\rangle=-5\cdot0.5=-2.5.
\]
Again FW direction is chosen:
\[
\gamma_1=0.1,\qquad x_2=0.5+0.1\cdot4.5=0.95.
\]
Objective: $f(x_2)=(0.95-3)^2\approx4.2025$.

\item \textbf{Iteration 2}
\begin{align*}
g_2&=2(0.95-3)=-4.1,\\
s_2&=5,\qquad v_2=\argmax_{v\in\{0,5\}}(-4.1 v)=0,\\
d^{\text{FW}}_2&=5-0.95=4.05,\qquad d^{\text{A}}_2=0.95-0=0.95,\\
\langle g_2,d^{\text{FW}}_2\rangle&=-4.1\cdot4.05\approx-16.605,\\
\langle g_2,d^{\text{A}}_2\rangle&=-4.1\cdot0.95\approx-3.895.
\end{align*}
FW step again:
\[
\gamma_2=0.1,\qquad x_3=0.95+0.1\cdot4.05=1.355.
\]
Objective: $f(x_3)=(1.355-3)^2\approx2.706.
\]

The iterates move monotonically towards the optimum $w^\star=3$ while staying inside the polytope $[0,5]$.
\end{enumerate}

\newpage
\section*{Assumptions}
\begin{assumption}[Convexity and Smoothness]\label{as:convex}
$f:\mathbb{R}^d\rightarrow\mathbb{R}$ is convex and $L$‑smooth, i.e.
\[
\|\nabla f(x)-\nabla f(y)\|\le L\|x-y\|,\qquad\forall x,y\in\mathbb{R}^d.
\tag{A1}
\]
\end{assumption}
\begin{assumption}[Polytope Domain]\label{as:polytope}
\C=\conv(\V)$ with a finite vertex set $\V$; let $\diam(\C)=\max_{u,v\in\C}\|u-v\|$.
\tag{A2}
\end{assumption}
\begin{assumption}[Pyramidal Width]\label{as:pyramidal}
Define the \emph{pyramidal width} of $\C$ as
\[
\pyr:=\min_{\substack{x\in\C,\;d\neq 0\\
\A\subseteq\V,\;x\in\conv(\A)}}
\max_{v\in\A}\frac{\langle d,\,x-v\rangle}{\|d\|\,\diam(\C)}\;>0.
\tag{A3}
\]
\end{assumption}
Assumption (A3) holds for any polytope and quantifies how “pointed’’ the feasible region is; it is crucial for linear convergence of AFW.

\section*{Main Convergence Theorem}
\begin{theorem}[Linear convergence of AFW]\label{thm:linear}
Let Assumptions \ref{as:convex}–\ref{as:pyramidal} hold and denote the optimal value $f^\star:=\min_{x\in\C}f(x)$.  
Define the primal gap $h_k:=f(x_k)-f^\star$ and the duality gap
\[
\gap_k:=\langle \nabla f(x_k),\,x_k-s_k\rangle.
\]
Then for all iterations $k\ge 0$ the AFW iterates satisfy
\[
h_{k+1}\;\le\;\bigl(1-\rho\bigr)h_k,
\qquad\text{with}\qquad
\rho:=\frac{\pyr^{\,2}}{4L\,\diam(\C)^{2}} \in (0,1).
\tag{7}
\]
Consequently,
\[
h_k\;\le\;(1-\rho)^{k}\,h_0,
\qquad
\gap_k\;\le\;2L\,\diam(\C)^{2}(1-\rho)^{k}.
\tag{8}
\]
\end{theorem}
\begin{theorem}[Sublinear $O(1/k)$ rate without pyramidal width]\label{thm:sublinear}
If only (A1)–(A2) are assumed, AFW enjoys the standard Frank‑Wolfe $O(1/k)$ guarantee:
\[
h_k\;\le\;\frac{2L\,\diam(\C)^{2}}{k+2}.
\tag{9}
\]
\end{theorem}

\section*{Theoretical Properties}
\begin{itemize}[leftmargin=*]
\item \textbf{Stability.} The line‑search (5) guarantees that the step size $\gamma_k$ never exceeds the feasible bound (4), thus iterates remain in $\C$ for all $k$.
\item \textbf{Hyper‑parameter sensitivity.} AFW does not require an external learning‑rate schedule; the only “parameter’’ is the choice of line‑search (exact or back‑tracking). The linear rate constant $\rho$ depends solely on problem geometry ($L$, $\diam(\C)$, $\pyr$) and not on any user‑chosen step‑size.
\item \textbf{Comparison to baseline FW.} Classical FW enjoys only the $O(1/k)$ bound (9). The addition of away steps together with a positive pyramidal width yields the geometric reduction (7). In the worst case where $\pyr=0$ (e.g. a smooth curved domain) AFW reduces to the baseline.
\end{itemize}

\section*{Discussion}
The theorem shows that for polyhedral feasible sets AFW attains a linear convergence rate comparable to projected gradient methods, while preserving the projection‑free nature of Frank‑Wolfe. The constant $\rho$ is explicit; improving the geometry of $\C$ (e.g., by scaling or rotating) can increase $\pyr$ and thus accelerate convergence.

\newpage
\section*{Preliminaries (Definitions and Lemmas)}
\begin{lemma}[Descent Lemma]\label{lem:descent}
If $f$ is $L$‑smooth then for any $x,d\in\mathbb{R}^d$ and $\gamma\in[0,1]$,
\[
f(x+\gamma d)\le f(x)+\gamma\langle\nabla f(x),d\rangle+\frac{L}{2}\gamma^{2}\|d\|^{2}.
\tag{L1}
\]
\end{lemma}
\begin{proof}
Standard from $L$‑smoothness; see e.g. Nesterov (2004). \qedhere
\end{proof}

\begin{lemma}[Duality gap lower bounds primal gap]\label{lem:gap}
For any $x\in\C$,
\[
\gap(x):=\langle\nabla f(x),x-s(x)\rangle\;\ge\;h(x):=f(x)-f^\star.
\tag{L2}
\]
\end{lemma}
\begin{proof}
Convexity of $f$ yields
\[
f^\star\ge f(s(x))\ge f(x)+\langle\nabla f(x),s(x)-x\rangle,
\]
rearranging gives $f(x)-f^\star\le\langle\nabla f(x),x-s(x)\rangle$. \qedhere
\end{proof}

\begin{lemma}[Progress bound for a FW step]\label{lem:fw-progress}
If $d_k=d^{\text{FW}}_k$ then, using the exact line‑search (5),
\[
h_{k+1}\;\le\;h_k-\frac{\gap_k^{2}}{2L\diam(\C)^{2}}.
\tag{L3}
\]
\end{lemma}
\begin{proof}
Let $d:=d^{\text{FW}}_k=s_k-x_k$. By Lemma~\ref{lem:descent} with $\gamma\in[0,1]$,
\[
f(x_k+\gamma d)\le f(x_k)+\gamma\langle g_k,d\rangle+\frac{L}{2}\gamma^{2}\|d\|^{2}.
\]
Since $\|d\|\le\diam(\C)$ and $\langle g_k,d\rangle=-\gap_k$, the quadratic in $\gamma$ is minimized at
\[
\gamma^\star=\min\Bigl\{1,\;\frac{\gap_k}{L\|d\|^{2}}\Bigr\}\;.
\]
If $\gamma^\star=1$ then $h_{k+1}\le h_k-\gap_k+\frac{L}{2}\diam(\C)^{2}\le h_k-\frac{\gap_k^{2}}{2L\diam(\C)^{2}}$ because $\gap_k\le L\diam(\C)^{2}$.  
If $\gamma^\star<1$, plugging $\gamma^\star$ gives the same bound after simple algebra. \qedhere
\end{proof}

\begin{lemma}[Progress bound for an away step]\label{lem:away-progress}
If $d_k=d^{\text{A}}_k$ then,
\[
h_{k+1}\;\le\;h_k-\frac{\pyr^{\,2}}{2L\diam(\C)^{2}}\,\gap_k.
\tag{L4}
\]
\end{lemma}
\begin{proof}
For an away direction $d^{\text{A}}_k=x_k-v_k$, convexity of $f$ gives
\[
\langle g_k,d^{\text{A}}_k\rangle = \langle g_k,x_k-v_k\rangle
      = \langle g_k,x_k-s_k\rangle +\langle g_k,s_k-v_k\rangle
      = -\gap_k + \langle g_k,s_k-v_k\rangle .
\]
Since $v_k$ maximises $\langle g_k,v\rangle$ over the active set,
\[
\langle g_k,s_k-v_k\rangle\le\langle g_k,s_k\rangle-\langle g_k,v_k\rangle
      = -\gap_k -\bigl(\langle g_k,v_k\rangle-\langle g_k,s_k\rangle\bigr)
      \le -\gap_k .
\]
Thus $\langle g_k,d^{\text{A}}_k\rangle\le -2\gap_k$.
Applying Lemma~\ref{lem:descent} with $d:=d^{\text{A}}_k$ and using $\|d\|\le\diam(\C)$,
\[
f(x_k+\gamma d)\le f(x_k)-2\gamma\gap_k+\frac{L}{2}\gamma^{2}\diam(\C)^{2}.
\]
The maximal feasible $\gamma$ is $\gamma_{\max}= \frac{\alpha_{v_k}}{1-\alpha_{v_k}}\le \frac{1}{\pyr}\frac{\gap_k}{L\diam(\C)^{2}}$ (definition of pyramidal width, see Lacoste‑Julien \& Jaggi 2015). Optimising the quadratic yields
\[
h_{k+1}\le h_k-\frac{\pyr^{\,2}}{2L\diam(\C)^{2}}\gap_k .
\]
\qedhere
\end{proof}

\section*{Main Proof (Step‑by‑Step Derivation)}
We now combine Lemmas~\ref{lem:fw-progress} and \ref{lem:away-progress} to obtain the linear rate (7).  
All non‑trivial steps are numbered.

\begin{proof}[Proof of Theorem \ref{thm:linear}]
Let $h_k:=f(x_k)-f^\star$ and $\gap_k$ be the duality gap at iteration $k$.

\textbf{[Step 1]} By Lemma~\ref{lem:gap}, $\gap_k\ge h_k$ for all $k$.

\textbf{[Step 2]} Consider the two possible cases for the direction $d_k$.

\begin{enumerate}[label=(\alph*)]
\item \emph{FW step.} By Lemma~\ref{lem:fw-progress},
\[
h_{k+1}\le h_k-\frac{\gap_k^{2}}{2L\diam(\C)^{2}}.
\tag{S2a}
\]
Using $\gap_k\ge h_k$ (Step 1) we obtain
\[
h_{k+1}\le h_k\Bigl(1-\frac{h_k}{2L\diam(\C)^{2}}\Bigr)
\le \bigl(1-\rho\bigr)h_k,
\]
where we have defined
\[
\rho:=\frac{\pyr^{\,2}}{4L\diam(\C)^{2}}.
\tag{S2a‑rho}
\]
The inequality $h_k\ge \pyr^{\,2}/2$ follows from the definition of $\pyr$ and the fact that $h_k\le L\diam(\C)^{2}$ (smoothness bound). Hence (S2a) implies $h_{k+1}\le(1-\rho)h_k$.

\item \emph{Away step.} By Lemma~\ref{lem:away-progress},
\[
h_{k+1}\le h_k-\frac{\pyr^{\,2}}{2L\diam(\C)^{2}}\gap_k.
\tag{S2b}
\]
Again applying $\gap_k\ge h_k$ (Step 1) yields
\[
h_{k+1}\le h_k\Bigl(1-\frac{\pyr^{\,2}}{2L\diam(\C)^{2}}\Bigr)
= (1-2\rho)h_k\le(1-\rho)h_k,
\tag{S2b‑final}
\]
since $0<\rho\le 1/2$ by definition.

\end{enumerate}
Thus in both cases we have the uniform bound
\[
h_{k+1}\le (1-\rho)h_k.
\tag{S3}
\]
\textbf{[Step 3]} Recursively applying (S3) yields
\[
h_k\le (1-\rho)^{k}h_0,
\tag{S4}
\]
which is exactly (7).  

\textbf{[Step 4]} To bound the duality gap, use Lemma~\ref{lem:gap} and the smoothness inequality
\[
\gap_k\le \|\nabla f(x_k)\|\,\diam(\C)\le L\diam(\C)^{2},
\]
together with $h_k\le (1-\rho)^{k}h_0\le (1-\rho)^{k}L\diam(\C)^{2}$. Combining with $\gap_k\le 2L\diam(\C)^{2}(1-\rho)^{k}$ gives (8). \qedhere
\end{proof}

\section*{Rate Derivation (With Constants)}
From Theorem~\ref{thm:linear} we have the explicit linear contraction factor
\[
\rho = \frac{\pyr^{\,2}}{4L\,\diam(\C)^{2}}.
\]
Hence after $k$ iterations
\[
f(x_k)-f^\star \le \bigl(1-\tfrac{\pyr^{\,2}}{4L\diam(\C)^{2}}\bigr)^{k}\,(f(x_0)-f^\star).
\]
If one prefers a more convenient bound, using $1-\alpha\le e^{-\alpha}$ gives
\[
f(x_k)-f^\star \le \exp\!\Bigl(-\frac{\pyr^{\,2}}{4L\diam(\C)^{2}}\,k\Bigr)\,(f(x_0)-f^\star).
\]

\section*{Special Cases}
\begin{itemize}[leftmargin=*]
\item \textbf{Strongly convex $f$.} If $f$ is $\mu$‑strongly convex, the primal gap dominates the squared distance to the optimum, and the same $\rho$ yields a linear rate on $\|x_k-x^\star\|^{2}$.
\item \textbf{Degenerate polytope (e.g. simplex).} For the probability simplex $\Delta_{d}$, $\pyr=1/\sqrt{d}$, $\diam(\Delta_{d})=\sqrt{2}$, so $\rho = \frac{1}{8d\,L}$, matching known rates for AFW on the simplex.
\item \textbf{Absence of away steps.} Setting $\pyr=0$ collapses $\rho$ to $0$; the bound (7) reduces to the classical $O(1/k)$ rate (9), recovering the standard Frank‑Wolfe guarantee.
\end{itemize}

\end{document}